\subsection{Характеристика регионов 1--5}

IF97 выделяет пять основных регионов и отдельную область насыщения.
Регионы различаются физическим смыслом
и используемым термодинамическим потенциалом.
Сводная характеристика приведена в таблице~\ref{tab:if97_regions_overview}.

\begin{table}[htbp]
\centering
\caption{Регионы IAPWS-IF97 и используемые формулировки}
\label{tab:if97_regions_overview}
\begin{tblr}{
  width=\linewidth,
  colspec={X[1,l] X[2,l] X[2,l]},
  row{1}={font=\bfseries},
  hline{1,2,Z} = {solid},
}
Регион & Физический смысл & Потенциал и независимые переменные \\
1 & Сжатая вода (подохлаженная жидкость) & Гиббс, $(p, T)$ \\
2 & Перегретый пар & Гиббс, $(p, T)$ \\
3 & Высокоплотная околокритическая область & Гельмгольц, $(\rho, T)$ \\
4 & Линия насыщения и двухфазная область & Насыщение, $(p, x)$ или $(T, x)$ \\
5 & Высокотемпературная область (пар) & Гиббс, $(p, T)$ \\
\end{tblr}
\end{table}

В инженерных сценариях используются входные пары
$(p, T)$, $(p, h)$, $(p, s)$, $(p, x)$ и $(\rho, T)$.
На выходе формируется унифицированное состояние:
$p$, $T$, $v$, $\rho$, $h$, $s$, $u$, $c_p$, $w$ и номер региона IF97.

\paragraph*{Регион 1: сжатая вода.}
Регион~1 описывает стабильную однофазную жидкость
	в диапазоне температур от $273{,}15$ до $623{,}15$~К
	при давлениях от линии насыщения до $100$~МПа~\cite{iapws_if97_2012}.
Для этого региона используется формулировка через потенциал Гиббса,
	а безразмерные переменные нормируются константами
	$p^{*}=16{,}53$~МПа и $T^{*}=1386$~К.
Прямая аппроксимация задаётся суммой из $34$ табличных коэффициентов
$n_i$ и степеней $I_i$, $J_i$.

	Особенностью региона~1 является согласование
с термодинамическим нулём на тройной точке:
коэффициенты уравнения подобраны так,
чтобы удельные внутренняя энергия и энтропия
насыщенной жидкости обращались в нуль.
Стандарт также допускает физически корректное описание
метастабильной жидкости в узкой области,
примыкающей к линии насыщения.

	Для инженерных обратных расчётов
	в регионе~1 предусмотрены аналитические зависимости
		$T(p, h)$ и $T(p, s)$.
Обе зависимости также задаются табличными наборами коэффициентов
по $20$ коэффициентов в каждой аппроксимации.
Это позволяет избегать итераций
в основной жидкофазной части рабочей области.

\paragraph*{Регион 2: перегретый и метастабильный пар.}
Регион~2 описывает перегретый пар
	и газообразное состояние воды
	в трёх температурно-барических диапазонах:
	при $273{,}15 \le T \le 623{,}15$~К и $0 < p \le p_{\mathrm{sat}}(T)$,
	при $623{,}15 < T \le 863{,}15$~К и $0 < p \le p_{B23}(T)$,
	а также при $863{,}15 < T \le 1073{,}15$~К и $0 < p \le 100$~МПа
	~\cite{iapws_if97_2012}.

	Математически регион~2 задаётся через энергию Гиббса,
	разбитую на идеально-газовую и остаточную части.
	В качестве нормировочных констант
	используются $p^{*}=1$~МПа и $T^{*}=540$~К.
Идеально-газовая часть аппроксимации
содержит $9$ табличных коэффициентов,
остаточная часть --- $43$ коэффициента.
	Отдельно стандарт учитывает метастабильный пар
	в примыкающей к линии насыщения области;
	для диапазона давлений до $10$~МПа
	используется специальная модификация остаточной части.

Для обратных расчётов
регион дополнительно разбивается на подобласти 2a, 2b и 2c.
Это разбиение используется только для аналитических зависимостей
$T(p, h)$ и $T(p, s)$,
которые задаются отдельными табличными наборами коэффициентов
для каждой подобласти,
и ускоряют расчёты в паровой области
без внешнего итерационного контура.

\paragraph*{Регион 3: высокоплотная околокритическая область.}
Регион~3 охватывает область
	от $623{,}15$ до $863{,}15$~К
	между границей $B23$ и верхним пределом давления $100$~МПа
	~\cite{iapws_if97_2012}.
	Именно здесь расположены критическая точка воды
	с параметрами $T_{c}=647{,}096$~К,
	$p_{c}=22{,}064$~МПа
	и $\rho_{c}=322$~кг/м\textsuperscript{3},
	а также наиболее выраженные нелинейности термодинамических свойств.

	В отличие от регионов~1, 2 и~5,
	для региона~3 независимыми переменными служат плотность и температура.
Поэтому стандарт использует не потенциал Гиббса,
а безразмерную энергию Гельмгольца,
нормированную критическими значениями
$\rho^{*}=\rho_{c}$ и $T^{*}=T_{c}$.
Такой выбор естественен для околокритической области,
где давление удобнее рассматривать
как производную по плотности,
а не как независимый аргумент.
	Базовая аппроксимация региона~3
	использует $40$ табличных коэффициентов
	$n_i$, $I_i$, $J_i$.

	Для региона~3 в IF97 не вводятся аналитические обратные уравнения
		вида $T(p, h)$ или $T(p, s)$.
Поэтому при постановках $(p, T)$, $(p, h)$ и $(p, s)$
требуется либо использование специальных вспомогательных релизов,
либо численное решение нелинейных уравнений.
Именно эта особенность определяет необходимость
дополнительных численных методов в проекте.

\paragraph*{Регион 4: линия насыщения и двухфазная область.}
Регион~4 описывает не двумерную область состояния,
	а кривую насыщения от тройной точки до критической точки:
	$273{,}15 \le T \le 647{,}096$~К
	и $611{,}213$~Па $\le p \le 22{,}064$~МПа
	~\cite{iapws_if97_2012}.
Для него не вводится собственное фундаментальное уравнение
в виде энергии Гиббса или Гельмгольца.
Вместо этого используются аналитически согласованные зависимости
$p_{\mathrm{sat}}(T)$ и $T_{\mathrm{sat}}(p)$,
получаемые из одного неявного полиномиального соотношения.
Это соотношение содержит $10$ табличных констант
$n_1,\dots,n_{10}$.

	Полные свойства на линии насыщения
	вычисляются через соседние однофазные регионы.
	Ниже $623{,}15$~К насыщенная жидкость
	и насыщенный пар вычисляются по формулам регионов~1 и~2,
	а в верхней околокритической части линии насыщения
	используется формулировка региона~3.
Внутри двухфазной области состояние дополнительно задаётся
степенью сухости $x$,
по которой выполняется расчёт свойств смеси.

\paragraph*{Регион 5: высокотемпературная область.}
Регион~5 используется для высокотемпературного пара
	при $1073{,}15 < T \le 2273{,}15$~К
	и давлениях до $50$~МПа~\cite{iapws_if97_2012}.
	Как и регион~2,
	он описывается через сумму идеально-газовой
	и остаточной частей энергии Гиббса.
	Для нормировки применяются
	$p^{*}=1$~МПа и $T^{*}=1000$~К.
Идеально-газовая и остаточная части
содержат по $6$ табличных коэффициентов.

	Для региона~5 в стандарте не предусмотрены
	аналитические обратные уравнения.
Поэтому при расчётах по энтальпии или энтропии
температура в этой области
должна восстанавливаться численными методами.
Кроме того, модель относится к химически чистой,
недиссоциированной воде
и не учитывает эффекты термической диссоциации,
которые при столь высоких температурах
могут становиться значимыми в специальных приложениях.
